\documentclass{article}

\usepackage[margin=0.85in]{geometry}
\usepackage{booktabs}
\usepackage{cite}
\usepackage{enumitem}
\usepackage{hyperref}

\setlist{nosep, topsep=2pt, leftmargin=1.5em}

\hypersetup{
    colorlinks=true,
    linkcolor=blue,
    citecolor=blue,
    urlcolor=blue
}

\title{\textbf{YZV411E - Big Data Analytics Project Proposal} \\ \large Distributed Analysis of Aircraft Turnaround Buffers and Delay Propagation Using Apache Spark}
\author{Aybek Taha At\i lgan \quad (ID: \dots)}
\date{October 2026}

\begin{document}

\maketitle

\section*{Abstract}
\textbf{Problem Definition:} Commercial flight delays rarely remain isolated at their point of origin. When an aircraft flies consecutive legs throughout a day, an inbound arrival delay either gets absorbed by the scheduled ground turnaround buffer or cascades into the subsequent departure as a reactionary delay. Existing academic work predominantly casts delay analysis as isolated binary classification on independent rows, ignoring rotational aircraft dependencies. Conversely, network queuing models in operations research capture delay trees but rely on single-machine solvers that cannot scale beyond single-airline slices or isolated days. To overcome these limitations, this project implements a distributed PySpark pipeline across 31.7 million U.S. flight records from 2018 through 2022 ($\approx 9.8$ GB uncompressed). By computing distributed chronological window functions over unique aircraft registrations (\texttt{Tail\_Number}), we quantify buffer absorption dynamics across major hubs and benchmark distributed execution against a single-node Pandas baseline.

\section{Project Aim and Motivation}
Aviation delays cost airlines and passengers billions of dollars annually. Secondary (reactionary) delays account for nearly half of total delay minutes in congested airspaces. When tight turnaround schedules fail to absorb an incoming disruption, the downstream departure is delayed despite clear local weather.

This study investigates three operational questions:
\begin{enumerate}
    \item What is the critical scheduled turnaround buffer ($B_{sched}$) required at major hubs to absorb incoming arrival delays without triggering downstream reactionary delays?
    \item How do hub airports differ in systemic resilience, distinguishing delay-absorbing hubs from delay-amplifying hubs across operating hours?
    \item What is the quantitative propagation multiplier ($R_0$) of root-cause delays across consecutive flight legs?
\end{enumerate}

\noindent\textbf{Why Big Data Technologies Are Required:} Tracking aircraft itineraries across five calendar years requires sorting and stateful windowing across 31.7 million records. A narrow-body aircraft executes 4 to 8 legs daily, forming continuous operational chains. Reconstructing these sequences requires a distributed stateful operation:
\begin{center}
\texttt{Window.partitionBy("Tail\_Number").orderBy("FlightDate", "CRSDepTime")}
\end{center}
On a single machine with 16 GB RAM, conventional libraries like Pandas crash with out-of-memory (OOM) errors or induce severe swap thrashing during global grouping and lagging. Apache Spark distributes airframe partitions across worker cores, manages disk spill gracefully, and utilizes Catalyst query optimization to prune unneeded columns prior to shuffle sorting.

\section{Data}
The project uses the \textit{Flight Delay Dataset 2018--2022}, sourced from the U.S. Bureau of Transportation Statistics (BTS) and hosted on Kaggle (\url{https://www.kaggle.com/datasets/robikscube/flight-delay-dataset-20182022}).
\begin{itemize}
    \item \textbf{Type and Volume:} Tabular relational flight operational logs across five yearly files (\texttt{Combined\_Flights\_2018.csv} through \texttt{2022.csv}), totaling \textbf{9.8 GB uncompressed CSV} ($\approx 31.7$ million rows, 61 features), comfortably exceeding the 1 GB course minimum. Compressed Parquet equivalents occupy $\approx 1.06$ GB.
    \item \textbf{Key Features:} Identifiers (\texttt{Tail\_Number}, \texttt{FlightDate}, \texttt{Operating\_Airline}), route endpoints (\texttt{Origin}, \texttt{Dest}), scheduled/actual timestamps (\texttt{CRSDepTime}, \texttt{DepTime}, \texttt{CRSArrTime}, \texttt{ArrTime}), delay durations (\texttt{DepDelay}, \texttt{ArrDelay}), and BTS delay causes (\texttt{LateAircraftDelay}, \texttt{CarrierDelay}, \texttt{WeatherDelay}, \texttt{NASDelay}).
    \item \textbf{Chunking and Partitioning Strategy:} Data is partitioned on disk by \texttt{Year} and \texttt{Month} in columnar Parquet format. During sequential linking, Spark DataFrames are hash-repartitioned by \texttt{Tail\_Number}. This ensures all chronological flights for an airframe reside within the same executor memory partition, eliminating redundant cross-network shuffles during lag calculations.
\end{itemize}

\section{Tools and Frameworks}
\begin{itemize}
    \item \textbf{Apache Spark (PySpark):} Primary distributed framework for ingestion, data hygiene, hash repartitioning, and windowed sequence transformations across multi-core executors.
    \item \textbf{Apache Parquet:} High-performance columnar storage with Snappy compression, shrinking storage from 9.8 GB to $\approx 1.1$ GB and enabling push-down predicates on \texttt{Origin} and \texttt{FlightDate}.
    \item \textbf{Pandas / DuckDB Baseline:} A single-node control benchmark on 500 MB, 2 GB, and 5 GB slices to document memory exhaustion thresholds and measure empirical speed-up ratios.
    \item \textbf{Plotly / Matplotlib:} Generation of turnaround response curves, hub propagation heatmaps, and scalability benchmark profiles.
\end{itemize}

\section{Expected Outcomes and Evaluation}
\begin{enumerate}
    \item \textbf{Benchmarked Distributed Performance:} Empirical comparison of single-node Pandas versus PySpark across scaled data volumes (500 MB to 9.8 GB). Metrics include wall-clock runtime, peak heap memory, CPU core utilization, and speed-up curves.
    \item \textbf{Turnaround Buffer Absorption Curves:} Quantification of net delay variation $\Delta D = \text{DepDelay}_{k+1} - \text{ArrDelay}_k$ against scheduled buffer $B_{sched} = \text{CRSDepTime}_{k+1} - \text{CRSArrTime}_k$, establishing the threshold where $\Delta D \le 0$ with 80\% confidence.
    \item \textbf{Hub Resilience Taxonomy:} Categorization of major hubs (ATL, ORD, DFW, DEN, JFK) into delay-absorbing versus delay-amplifying nodes, controlling for seasonal weather and airport congestion.
    \item \textbf{Reproducible Pipeline and Insights:} Fully documented PySpark codebase and interactive dashboard heuristics for buffer schedule planning.
\end{enumerate}

\section{Progress Plan}
The following milestones are planned prior to the progress report submission:
\begin{itemize}
    \item \textbf{Step 1 (Weeks 1--2) -- Ingestion and Quality Audit:} Download 2018--2022 files, compute SHA-256 checksums, audit \texttt{Tail\_Number} missingness, and build the partitioned Parquet conversion pipeline.
    \item \textbf{Step 2 (Weeks 3--4) -- Baseline Profiling and Validation:} Profile delay distributions across carriers/seasons and validate BTS \texttt{LateAircraftDelay} tags against ground-truth delays.
    \item \textbf{Step 3 (Weeks 5--6) -- Sequence Linking Pipeline:} Construct the PySpark windowing pipeline, handling boundary edge cases (overnight layovers $>8$ hours and diverted legs).
    \item \textbf{Step 4 (Weeks 7--8) -- Benchmarking and Propagation Modeling:} Execute standardized runtime/memory benchmarks (Pandas vs. PySpark) and compute hub buffer absorption curves.
    \item \textbf{Step 5 (Weeks 9--10) -- Visualization and Final Deliverable:} Produce interactive dashboards, synthesize operational recommendations, and author the final report.
\end{itemize}

\begin{thebibliography}{99}\small
\bibitem{beatty2001}
R.~Beatty et al., ``Evaluation of flight delay propagation through an airline schedule,'' in \textit{4th USA/Europe ATM R\&D Seminar}, 2001.
\bibitem{pyrgiotis2013}
N.~Pyrgiotis et al., ``A complex network approach to modeling delay propagation,'' \textit{Transp. Res. Part C}, vol.~27, pp.~208--238, 2013.
\bibitem{fleurquin2013}
P.~Fleurquin et al., ``Systemic delay propagation in the US airport network,'' \textit{Sci. Rep.}, vol.~3, p.~1159, 2013.
\bibitem{bts2023}
Bureau of Transportation Statistics, ``Reporting causes of flight delays,'' U.S. DOT, Tech. Rep., 2023.
\end{thebibliography}

\end{document}
